%----------
%--- CONCLUSOES

\chapter{CONSIDERAÇÕES FINAIS}
\thispagestyle{myheadings}

No início deste trabalho, verificou-se que os processos de ETL da SEFAZ-PB não acompanhavam 
o crescimento do volume de dados, comprometendo as atividades de fiscalização. Em resposta 
a esse cenário, investigou-se a viabilidade de uma reestruturação do ETL baseada em 
tecnologias de Big Data, visando ampliar o desempenho, a manutenibilidade e a eficiência 
no armazenamento dos dados.

Desse modo, o objetivo geral desta pesquisa, de reestruturar a estratégia de ETL com base no
projeto BDMalhas, foi alcançado. Essa conclusão apoia-se na comparação entre o sistema legado e a
nova implementação apresentada no \chapref{cap:resultados}: a \tabref{tab:comparacao_manutenibilidade}
evidenciou melhorias na estrutura dos fluxos, a \tabref{tab:comparacao_armazenamento} evidenciou a
redução no espaço em disco utilizado, e a \tabref{tab:desempenho_cargas} apontou resultados
favoráveis de desempenho em determinados cenários, embora sua avaliação comparativa tenha sido
limitada pelas diferenças entre os ambientes e volumes de dados processados.

Quanto ao primeiro objetivo específico, este foi alcançado por meio do desenvolvimento de
um sistema de ETL que materializou a proposta de reestruturação, utilizando Airflow para
orquestração, Spark para processamento distribuído e ClickHouse como banco de dados
analítico, conforme detalhado no \chapref{chap:desenvolvimento}.

Em seguida, a avaliação da nova implementação permitiu atender ao segundo objetivo específico.
O ganho de manutenibilidade foi verificado por meio da comparação do número de nós, de
dependências entre nós e de consultas SQL entre as duas implementações, apresentada na
\tabref{tab:comparacao_manutenibilidade}, que evidenciou a simplificação dos \textit{pipelines}
de ETL. 

A redução no uso de armazenamento, por sua vez, foi constatada a partir da comparação do
espaço ocupado pelas tabelas do sistema legado e da nova solução, apresentada na
\tabref{tab:comparacao_armazenamento}. 

No mais, os resultados favoráveis de desempenho, apresentados na \tabref{tab:desempenho_cargas} e
no Gráfico \ref{graf:tempo_exec}, reforçam o atendimento a esse objetivo. Essa comparação, no
entanto, não é direta, dadas as diferenças de volume e de ambiente entre as implementações,
conforme discutido na \secref{sec:limitacoes_comparacao}.

Por fim, o \chapref{chap:discussao} permitiu alcançar o terceiro objetivo específico. Nele,
identificam-se os principais desafios enfrentados durante a reestruturação e os ganhos obtidos em
termos de manutenibilidade e de economia de espaço de armazenamento, além de indicativos dos
possíveis ganhos de escalabilidade proporcionados pela nova arquitetura.

Com isso, os resultados obtidos fornecem evidências favoráveis à hipótese de que a 
reestruturação dos \textit{pipelines} de ETL pode trazer impactos positivos para os 
processos de Big Data da SEFAZ-PB, indicando a viabilidade da abordagem proposta.

Diante dos resultados obtidos, conclui-se que o problema de pesquisa foi respondido de forma
satisfatória nas dimensões avaliadas. A \tabref{tab:comparacao_manutenibilidade} e a
\tabref{tab:comparacao_armazenamento} confirmam ganhos conclusivos de manutenibilidade e de
armazenamento, respectivamente. Já os resultados de desempenho, apresentados na
\tabref{tab:desempenho_cargas}, foram favoráveis em determinados cenários, mas não permitem uma
comparação definitiva entre as duas implementações, conforme discutido na
\secref{sec:limitacoes_comparacao}. Ainda assim, a pesquisa demonstra que a reestruturação da
estratégia de ETL baseada em tecnologias de Big Data constitui uma alternativa viável para
modernizar os processos da SEFAZ-PB.

Quanto às contribuições desta pesquisa, destacam-se três frentes. No âmbito prático, entrega-se à
SEFAZ-PB um sistema de ETL funcional para os quatro fluxos avaliados, que demonstra potencial para
substituir a solução legada baseada em SSIS, ainda que sua adoção em produção dependa da conclusão
das adequações do \textit{Data Lake} discutidas ao longo deste trabalho.

No âmbito técnico, destaca-se a abstração \textit{BaseTask}, desenvolvida para padronizar a criação e a configuração
dos operadores do Airflow, que se consolidou como um padrão reaproveitável para a construção de
\textit{pipelines} com estrutura repetitiva, constituindo uma das principais contribuições de
engenharia deste trabalho. Por fim, no âmbito do conhecimento, esta pesquisa documenta um caso real
de migração de um sistema de ETL de SSIS para Airflow e Spark em um órgão público, cenário pouco
explorado na literatura sobre reestruturação de processos de ETL.

Metodologicamente, esta pesquisa foi conduzida como um estudo de caso no projeto BDMalhas,
adotando uma abordagem quantitativa para desenvolver e avaliar a proposta de reestruturação do ETL. A avaliação 
da implementação baseou-se em experimentos comparativos realizados no ambiente 
disponibilizado pela SEFAZ-PB.

Entretanto, por tratar-se de um estudo de caso, os resultados obtidos estão limitados ao contexto
do projeto BDMalhas e ao ambiente experimental utilizado. Além disso, as avaliações foram 
realizadas sobre uma cópia do \textit{Data Lake} contendo apenas dados compreendidos entre 
janeiro e março de 2020, com volume inferior ao disponível no ambiente de produção. Dessa forma, 
os experimentos ficaram restritos a esse período e contexto, limitando a generalização dos 
resultados para outros cenários.

Por fim, como recomendação para trabalhos futuros, sugere-se a avaliação da proposta em um 
ambiente adequado, com infraestrutura e volumes de dados mais próximos aos encontrados em 
produção, de modo a permitir comparações mais conclusivas de desempenho entre as duas 
soluções. Recomenda-se também a realização de novos experimentos com diferentes períodos e 
volumes de dados.

Além disso, sugere-se a utilização de um ambiente com maior capacidade computacional e um 
número maior de \textit{workers} no cluster Spark, permitindo avaliar de forma mais abrangente 
os benefícios do processamento distribuído e sua escalabilidade. Recomenda-se, ainda, a avaliação 
da arquitetura em outros processos de ETL da SEFAZ-PB e a investigação de estratégias de 
otimização que possam ampliar sua escalabilidade, desempenho e manutenibilidade.
